% ATS_Repair_Rigidity_v0.3_de.tex — Gap-Serie Teil 6
% R18-Korrekturfassung, 2. Oktober 2026. Quellengebundene Anspruchsgrenzen.
\input{glyphtounicode}
\pdfgentounicode=1
\documentclass[11pt,a4paper]{article}
\usepackage{cmap}
\usepackage[utf8]{inputenc}
\usepackage[T1]{fontenc}
\usepackage[ngerman]{babel}
\usepackage{lmodern}
\usepackage{microtype}
\emergencystretch=3em
\usepackage{amsmath,amssymb,amsthm}
\usepackage{booktabs}
\usepackage{tabularx}
\usepackage[colorlinks=true,linkcolor=blue,citecolor=blue,urlcolor=blue]{hyperref}
\usepackage{xurl}

\renewcommand*{\HyperDestNameFilter}[1]{ger.#1}

\hypersetup{
  pdftitle={Was die Produktformel zulässt: Ein Starrheitssatz und ein Effektivgewichts-Test für die ATS-Höhen-Teilkette},
  pdfauthor={Lukas Geiger},
  pdfsubject={Zahlentheorie, Arithmetische Geometrie, Teichmüller-Theorie, abc-Vermutung},
  pdfkeywords={ATS, Arithmetischer Teichmüller-Raum, Produktformel, Starrheit, Projektive Höhe, Effektive Gewichte, Joshi},
  pdfcreator={LaTeX with hyperref},
  pdfproducer={pdfTeX},
  bookmarksnumbered=true,
  pdfpagemode=UseOutlines,
  pdfstartview={FitH},
  unicode=true
}

\newtheorem{theorem}{Satz}[section]
\newtheorem{lemma}[theorem]{Lemma}
\newtheorem{proposition}[theorem]{Proposition}
\newtheorem{corollary}[theorem]{Korollar}
\theoremstyle{definition}
\newtheorem{definition}[theorem]{Definition}
\theoremstyle{remark}
\newtheorem{remark}[theorem]{Bemerkung}

\newcommand{\OL}{\mathcal{O}_L}
\newcommand{\Geff}{\mathcal{G}_{\mathrm{eff}}}

\title{Was die Produktformel zulässt:\\Ein Starrheitssatz und ein Effektivgewichts-Test\\für die ATS-Höhen-Teilkette}
\author{Lukas Geiger\\\small Unabhängiger Forscher, Bernau, Deutschland\\\small ORCID: 0009-0005-7296-1534}
\date{2. Oktober 2026 --- Version 0.3 (Entwurf)}

\begin{document}
\maketitle

\begin{abstract}
Teil~5 dieser Serie~\cite{part5} berichtet ein quellengebundenes ATS-Substanzaudit;
die korrigierte Fassung von Teil~6 benennt die genaue Reichweite ihres
Effektivgewichts-Tests. Für einen festen Zahlkörper $L$, seine festen
Standard-Absolutbeträge und die volle Produktformel mit Tests an jedem
$x\in L^\times$ muss eine Familie reeller Exponenten $(\gamma_v)$ über alle
Stellen konstant sein, unabhängig von der Klassenzahl. Umverteilung mit
endlichen Fasern liefert dieselbe Folgerung für effektive Gewichte. Eine
vorgegebene $j^2$-Gewichtung auf einer nichtleeren Stellenmenge besitzt die
konstante globale Fortsetzung; ausgeschlossen sind nur ungleiche endgültige
Stellengewichte. Für globale Größenfamilien mit nachgewiesener linearer
Effektivgewichts-Darstellung lässt die Produktformel höchstens einen
$y$-abhängigen globalen Skalar zu. Bloße Abhängigkeit von den Gewichten genügt
nicht für Klassenmitgliedschaft. Unter ATS~II($\tfrac12$) (5.3.3) identifiziert
direktes Einsetzen den gewichteten Grad auf $L^\times$ und die projektive Höhe
auf $\mathbb P^n(L)$ mit ihren klassischen Werten; die Stabilisierung der
Quelle auf $L^\times$ ist dann summandenweise konstant. Die Erweiterung auf
$\mathbb P^n(\overline L)$, die ATS-III-Linebundle-Höhe und der
ATS-IV-Volumenapparat benötigen getrennte Brücken. Mitgliedschaft oder
Nichtmitgliedschaft von Log-Volumina und Maßen in der linearen Klasse wird
hier nicht klassifiziert. Koordinatenhyperebenen und ihre Rückbilder sind
verschiedene Objekte. Weder ein globaler Abschluss des ATS-Reparaturprogramms
noch ein Beweis oder eine Widerlegung von abc/IUT oder eine mathematische
Zertifizierung der gesamten ATS-Serie wird behauptet.
\end{abstract}

\clearpage
\tableofcontents
\clearpage

%=====================================================================
\section{Auftrag, Reichweite und Quellenversionen}\label{sec:mandate}

Dieser Beitrag ist Teil~6 der abc-Gap-Serie. Teil~1~\cite{part1} lokalisierte
die quantitative Lücke im Transportschritt von Korollar~3.12 aus IUTchIII~\cite{mochizukiIUT}; Teil~2~\cite{part2} lieferte ein Formkriterium; Teil~3~\cite{part3} destillierte einen prüfbaren
Zulässigkeitsstandard; Teil~4~\cite{part4} auditierte die ATS-Serie~\cite{joshiI,joshiII,joshiIIhalf,joshiIII,joshiIV} auf Ledger-Ebene (v1-F). Der
vorliegende Beitrag berichtet den mathematischen Kern des Reparaturprogramms:
einen Starrheitssatz, ein Faktorisierungsinstrument und deren
quellengebundene Anwendung auf die \emph{Höhen-Teilkette} der ATS-Texte.

Reichweitendisziplin. \glqq Die auditierten Passagen\grqq{} bezeichnet: ATS~II($\tfrac12$) \S\S5.3, 5.10, 8.2--8.7 und ATS~IV \S1.5, in genau den in der
Bibliographie zitierten arXiv-Versionen (II($\tfrac12$): v12; IV: v2; beide
revidiert am 24. Februar 2025); ATS~III v4, Remark~8.9.2(2), ist ein
gesondert typisiertes Anwendungsziel. Es handelt sich um öffentlich hinterlegte,
versionierte Preprints, nicht um begutachtete Publikationen; ATS~IV trägt die
Kopfzeile \glqq Preliminary version for comments\grqq{}. Dieser Beitrag auditiert nicht
die vollständige ATS-Serie, klassifiziert deren Volumen- und Maßapparat nicht
(siehe Fall C in Abschnitt~\ref{sec:consequences}) und erhebt keinen Anspruch
über Strukturen außerhalb der auditierten Passagen. Grundregeln wie in den
Teilen~4--5: kein Autoritätsargument, kein Wahrheitsverdikt, wörtliche Anker;
Befunde zu späterer Programmarbeit erscheinen nur als datierte externe
Statushinweise (Abschnitt~\ref{sec:notices}), getrennt von den Beweisen dieses
Beitrags.

%=====================================================================
\section{Typen und operatives Glossar}\label{sec:types}

Fixiert seien ein Zahlkörper $L$ und seine volle Menge $V_L$ archimedischer
und endlicher Stellen. Lokale Eingaben und globale Größenfamilien haben
verschiedene Typen:
\begin{itemize}
\item $b_v$ ist der feste Standard-Absolutbetrag: $|\sigma(x)|$ an einer
reellen Stelle, $|\sigma(x)|^2$ an einer komplexen Stelle und
$N\mathfrak p^{-v_{\mathfrak p}(x)}$ an einer endlichen Stelle. Daher gilt
$\prod_v b_v(x)=1$ für jedes $x\in L^\times$.
\item $\rho_{y,v}=b_v^{\,a_{y,v}}$ mit $a_{y,v}>0$ ist der äquivalente
Quellen-Absolutbetrag auf $L_v$. Das Quellendatum wird mit $y$ abgekürzt.
\item $\lambda_{y,v}$ ist ein Normalisierungsexponent. In unserer Notation
besagt ATS~II($\tfrac12$) (5.3.3) \cite{joshiIIhalf}
$b_v=\rho_{y,v}^{\,\lambda_{y,v}}$. Dies ist eine Paraphrase mit umbenannten
Symbolen, kein wörtliches Zitat der $\alpha/\xi$-Notation der Quelle.
\item $\beta_{y,v}=a_{y,v}\lambda_{y,v}$ ist das effektive Gewicht.
Die punktweise Identität (5.3.3) impliziert $\beta_{y,v}=1$ für jedes $v$
und außerdem $\lambda_{y,v}>0$.
\end{itemize}
Der Starrheitssatz beschränkt effektive Gewichte, nicht die beiden Faktoren
getrennt. Koordinatenvariation muss daher metrische Werte nicht verändern.

Für einen präzisen Koordinatenvergleich verwenden wir den gemeinsamen Raum
$\mathbb R^{(V_L)}$ endlich getragener Vektoren und
\[
H_y=\{t\in\mathbb R^{(V_L)}:\sum_v\lambda_{y,v}t_v=0\},
\qquad \ell_y(x)=(\log\rho_{y,v}(x))_v.
\]
Für $x\in L^\times$ hat $\ell_y(x)$ endlichen Träger.
Die Normalisierungsidentität liefert $\ell_y^{-1}(H_y)=L^\times$.
Nichtproportionale Koeffizientenfamilien definieren verschiedene Kerne in
diesem gemeinsamen Raum, während diese Rückzüge übereinstimmen.
Diese bedingte linear-algebraische Beobachtung beweist weder
Nichtproportionalität für ein konkretes Quellenparameterpaar noch bestätigt
sie Theorem~5.10.1 insgesamt.

Eine Familie \emph{faktorisiert über effektive Gewichte} hier nur mit der
linearen Darstellung aus Definition~\ref{def:geff}.
Ein \emph{klassischer} Wert wird aus der festen Standardfamilie auf dem
angegebenen arithmetischen Bereich berechnet. Eine
\emph{stellenselektive} Gewichtsfamilie hat ungleiche Endgewichte statt
eines einzigen globalen Skalars.

\begin{table}[htbp]
\centering\small
\caption{Lokale Eingaben, Koordinaten und globale Größenfamilien.}\label{tab:type_glossary}
\begin{tabularx}{\textwidth}{@{}lXX@{}}
\toprule
Objekt & Typ und Abhängigkeit & Bezug zu $\Geff$\\
\midrule
$b_v$ & Feste lokale Referenz & Eingabe, keine globale Familie\\
$\rho_{y,v}$ & Möglicherweise variierender lokaler Absolutbetrag & Eingabe\\
$\lambda_{y,v}$ & Normalisierungskoordinate & Parameter\\
$\beta_{y,v}$ & Effektiver Koeffizient; unter (5.3.3) identisch $1$ & Koeffizient\\
$H_y$ & Koordinatenkern; Variation bedingt durch Koeffizienten & Koordinatenobjekt\\
$h_y,\deg_y$ & Globale Werte auf $\mathbb P^n(L)$ und $L^\times$; unter (5.3.3) invariant & Geprüfte lineare Darstellungen\\
$\mathrm{vol}_y,\mu_y$ & Konkrete Maßobjekte und Abhängigkeit erfordern Identifikation & Zugehörigkeit und Nichtzugehörigkeit unklassifiziert\\
\bottomrule
\end{tabularx}
\end{table}

%=====================================================================
\section{Der Starrheitssatz}\label{sec:rigidity}

\begin{theorem}[Starrheit von Produktformel-Gewichten]\label{thm:rigidity}
Sei $L$ ein Zahlkörper und $(\gamma_v)_{v \in V_L}$ eine Familie reeller
Exponenten mit
\[
\prod_{v \in V_L} b_v(x)^{\gamma_v} = 1
\qquad \text{für alle } x \in L^\times .
\]
Dann ist $\gamma_v$ konstant in $v$. (Wir schreiben $\gamma$ für die
generische Exponentenfamilie des Satzes und reservieren $\lambda$ für die
Normalisierungsexponenten der Quelle aus Abschnitt~\ref{sec:types}.)
\end{theorem}

\begin{proof}[Beweis]
Setze $\mu_v := \gamma_v - 1$. Subtrahiert man den Logarithmus der
Standard-Produktformel vom Logarithmus der Voraussetzung, so ergibt sich
\[
\sum_{v \in V_L} \mu_v \log b_v(x) = 0
\qquad \text{für alle } x \in L^\times,
\]
wobei für jedes feste $x$ alle bis auf endlich viele Terme verschwinden.

\emph{(i) Archimedische Komponente.} Für Einheiten $u \in \OL^\times$
verschwinden alle endlichen Terme, also gilt
$\sum_{\sigma\,\mathrm{arch}} \mu_\sigma \log b_\sigma(u) = 0$ für alle $u$.
Nach dem Dirichletschen Einheitensatz~\cite[Ch.~I, \S7]{neukirch} ist
$\mathrm{Log}(\OL^\times)$ ein volles Gitter in der Spur-Null-Hyperebene
$H = \{ y \in \mathbb{R}^{r_1 + r_2} : \sum_\sigma y_\sigma = 0 \}$: Es liegt
in $H$, weil $\sum_\sigma \log b_\sigma(u) = \log|N(u)| = 0$, und sein Rang
ist $r_1 + r_2 - 1 = \dim H$. Der Vektor $(\mu_\sigma)_\sigma$ annulliert ein
volles Gitter von $H$, mithin ganz $H$, und liegt daher in
$H^\perp = \mathbb{R}\cdot(1,\dots,1)$:
\[
\mu_\sigma = c \quad \text{für alle archimedischen } \sigma,
\text{ für ein } c \in \mathbb{R}.
\]

\emph{(ii) Endliche Komponente.} Sei $\mathfrak{p}$ ein beliebiges Primideal
und $h$ die Ordnung seiner Klasse in der endlichen Klassengruppe~\cite[Ch.~I, \S6]{neukirch}, sodass $\mathfrak{p}^h = (x_{\mathfrak{p}})$ ein
Hauptideal ist. Dann ist $\log b_{\mathfrak{q}}(x_{\mathfrak{p}}) = 0$ für
alle endlichen $\mathfrak{q} \ne \mathfrak{p}$, während
$\log b_{\mathfrak{p}}(x_{\mathfrak{p}}) = -h \log N\mathfrak{p}$ und
\[
\sum_{\sigma} \log b_\sigma(x_{\mathfrak{p}})
= \log|N(x_{\mathfrak{p}})| = \log\bigl(N\mathfrak{p}^{\,h}\bigr)
= h \log N\mathfrak{p} .
\]
Einsetzen in die Annullierungsidentität und Verwendung von (i) liefert:
\[
\mu_{\mathfrak{p}} \cdot (-h \log N\mathfrak{p})
+ c \cdot (h \log N\mathfrak{p}) = 0
\quad\Longrightarrow\quad
\mu_{\mathfrak{p}} = c .
\]
Da $\mathfrak{p}$ beliebig war, gilt $\mu_v = c$ für alle $v$, d.\,h.\
$\gamma_v = 1 + c$ ist konstant. Umgekehrt erfüllt jede konstante Familie die
Voraussetzung.
\end{proof}

\begin{remark}[Status]\label{rem:status}
Der Satz ist elementar und liegt im Umfeld der Artin--Whaples-Axiomatik~\cite{artinwhaples}; wir formulieren und beweisen ihn vollständig, weil die
Serie ihn als typisiertes Prüfinstrument verwendet und weil die in den
Abschnitten~\ref{sec:application}--\ref{sec:consequences} gezogene Konsequenz
von seinen genauen Voraussetzungen abhängt. Für den Satz selbst wird kein
Neuheitsanspruch erhoben.
\end{remark}

\begin{corollary}[Ungleiche Endgewichte sind ausgeschlossen]\label{cor:noselect}
Unter der Voraussetzung von Satz~\ref{thm:rigidity} für denselben Körper
und alle Stellen erfüllen nur konstante Exponentenfamilien die Produktformel.
Insbesondere erzwingt die Vorgabe $\gamma_v=j^2$ auf einer nichtleeren Menge
$S\subseteq V_L$ den Wert $\gamma_v=j^2$ an jeder Stelle.
Die konstante globale Fortsetzung ist zulässig; eine Fortsetzung durch
ungleiche komplementäre Gewichte ist unmöglich.
\end{corollary}

\begin{lemma}[Umverteilung mit endlichen Fasern]\label{lem:redistribution}
Sei $\tau : V_L \to V_L$ eine Abbildung mit endlichen Fasern (insbesondere sei
$\tau$ eine Permutation), sei $a_v > 0$, und die Familie
$\bigl(b_{\tau(v)}^{\,a_v \gamma_v}\bigr)_{v}$ erfülle die Produktformel für
alle $x \in L^\times$. Dann sind die effektiven Gewichte
$\beta_u := \sum_{v :\, \tau(v) = u} a_v \gamma_v$ konstant in $u$.
\end{lemma}

\begin{proof}[Beweis]
Umindizierung der (für jedes feste $x$ endlichen) Summe liefert
$\sum_v a_v\gamma_v \log b_{\tau(v)}(x) = \sum_u \beta_u \log b_u(x) = 0$ für
alle $x \in L^\times$ --- das ist die Voraussetzung von
Satz~\ref{thm:rigidity} in den effektiven Gewichten. Die Endlichkeit der
Fasern macht jedes $\beta_u$ zu einer endlichen Summe und die Umindizierung
zulässig.
\end{proof}

Für unendliche Korrespondenzen bräuchte man zusätzlich absolute
Summierbarkeit je Faser und eine gerechtfertigte Umordnung; diese
Allgemeinheit wird in diesem Beitrag nicht verwendet, und für die auditierten
Passagen, die keine Stellen bewegen, wird sie auch nicht benötigt.

%=====================================================================
\section{Die Effektivgewichts-Größenklasse und die
Aufzehrungsproposition}\label{sec:application}

\begin{definition}[Effektivgewichts-Größenklasse]\label{def:geff}
Fixiert sei ein Bereich $Z$. Eine Familie $Q=(Q_y:Z\to\mathbb R)_y$
liegt in $\Geff$, wenn es $y$-unabhängige lokale Funktionale
$q_w:Z\to\mathbb R$, $w\in V_L$, gibt mit
\[
Q_y(z)=\sum_{w\in V_L}\beta_{y,w}q_w(z),
\qquad Q_{\mathrm{ref}}(z):=\sum_{w\in V_L}q_w(z).
\]
Beide Summen müssen endlich oder absolut konvergent sein, die erste für
jedes $(y,z)$ und die Referenzsumme für jedes $z$.
Die Anforderung an die Referenzsumme erfasst auch Parameter mit $\beta_y=0$.
\end{definition}

Zugehörigkeit erfordert diese \emph{lineare} Darstellung; bloße Abhängigkeit
von $\beta_y$ reicht nicht. Auf einem einpunktigen Bereich setze man etwa
$\beta_{y,w}=c_y$ an jeder Stelle, wobei $c_y$ die Werte $1$ und $2$
annimmt, und $Q_y=c_y^2$. Eine Darstellung ergäbe $Q_y=c_y S$ mit
$y$-unabhängigem $S=\sum_wq_w$ und erzwänge sowohl $S=1$ als auch $S=2$.
Dieses Gegenbeispiel zur früheren Zugehörigkeitsbeschreibung gehört zur
R18-Korrektur und ist kein neuer Befund über eine Quellenhöhe.
Die generische Referenz $Q_{\mathrm{ref}}$ muss keine klassische
arithmetische Größe sein; diese Bezeichnung verlangt eine gesondert
identifizierte Höhe oder einen Grad samt Bereich.

\begin{lemma}[Faktorisierung]\label{lem:factorization}
Es gelte für den festen Körper und alle Stellen die gewichtete Produktformel
\[
\prod_v\rho_{y,v}(x)^{\lambda_{y,v}}=1
\qquad\text{für jedes }x\in L^\times.
\]
Dann ist $\beta_{y,w}=c_y$ für alle $w$, und jedes $Q\in\Geff$ erfüllt
\[
Q_y=c_yQ_{\mathrm{ref}}.
\]
Der Skalar darf von $y$ abhängen. Gilt (5.3.3) punktweise, so ist
$\beta_y=1$ und $Q_y=Q_{\mathrm{ref}}$.
Abstrakte reelle Gewichte erlauben $c_y=0$; für die positiven
$a_{y,v}$ und positiven Normalisierungsexponenten der Quelle gilt $c_y>0$.
\end{lemma}

\begin{proof}[Beweis]
Für jedes $x\in L^\times$ wird die Voraussetzung zu
\[
0=\sum_v\lambda_{y,v}\log\rho_{y,v}(x)
=\sum_w\beta_{y,w}\log b_w(x).
\]
Satz~\ref{thm:rigidity} liefert $\beta_{y,w}=c_y$.
Einsetzen in Definition~\ref{def:geff} ergibt die Behauptung; die festgelegte
Konvergenz gewährleistet die Existenz der Referenzsumme auch für $c_y=0$.
\end{proof}

\begin{proposition}[Aufzehrung auf den angegebenen arithmetischen Bereichen]\label{prop:exhaustion}
Für die projektive Höhe aus ATS~II($\tfrac12$), Definition~8.2.1
\cite{joshiIIhalf}, auf $\mathbb P^n(L)$ und den gewichteten Grad auf
$L^\times$ verwenden wir die feste volle Standardfamilie und die
Normalisierungsidentität (5.3.3). Dann gilt
\begin{align*}
h_y([x_0:\dots:x_n])
&=\sum_v\lambda_{y,v}\log\max_i\rho_{y,v}(x_i)\\
&=h_{\mathrm{cl}}([x_0:\dots:x_n]),\\
\deg_y(x)&=\sum_v\lambda_{y,v}\log\rho_{y,v}(x)
=\deg_{\mathrm{cl}}(x).
\end{align*}
Die Gradgleichheit hat den Bereich $x\in L^\times$.
\end{proposition}

\begin{proof}[Beweis]
Auf einem von null verschiedenen Repräsentantentupel
$(x_0,\dots,x_n)\in L^{n+1}\setminus\{0\}$ liefern die Positivität von
$a_{y,v}$ und $\beta_{y,v}=1$
\begin{align*}
\sum_v\lambda_{y,v}\log\max_i\rho_{y,v}(x_i)
&=\sum_v\lambda_{y,v}a_{y,v}\log\max_i b_v(x_i)\\
&=\sum_v\log\max_i b_v(x_i).
\end{align*}
Für jedes Tupel haben diese Summen endlichen Träger. Ein Wechsel des
Repräsentanten um $u\in L^\times$ ändert die Standardsumme um
$\sum_v\log b_v(u)=0$; die gewichtete Summe hat dieselbe Eigenschaft durch
(5.3.4) und (5.3.3), entsprechend Lemma~8.2.3 der Quelle.
Damit steigt die Gleichheit zu $\mathbb P^n(L)$ ab.
Für den Grad entfällt das Maximum; die Rechnung gilt auf $L^\times$.
\end{proof}

Die lokalen Höhenfunktionale
$q_v(x_0,\dots,x_n)=\log\max_i b_v(x_i)$ sind Funktionen auf
Repräsentantentupeln, nicht einzeln wohldefinierte projektive Funktionen.
Sie liefern eine lineare Darstellung auf dem Tupelbereich.
Für eine Darstellung auf $\mathbb P^n(L)$ kann man zunächst für jeden Punkt
einen $y$-unabhängigen Repräsentanten fixieren; die \emph{summierte} Höhe
ist unter der Produktformel von dieser Wahl unabhängig.
Für den Grad gilt $q_v(x)=\log b_v(x)$ auf $L^\times$.

\begin{remark}[Grenze und Stabilisierung]\label{rem:boundary}
Dies betrifft die identifizierten Bereiche und die lineare Klasse, nicht
jedes Höhen- oder Maßobjekt der ATS-Serie.
In ATS~II($\tfrac12$) \S\S8.6--8.7 wird die stabilisierte Höhe auf
$x\in L^\times$ ausgewertet. Jeder Summand $h_{\alpha\cdot y}(x)$ ist
unter der Normalisierungsidentität gleich demselben klassischen Wert;
sein Supremum ist daher ebenfalls dieser Wert. Für diese konstante Familie
ist kein zusätzlicher Tausch von Supremum und Summe nötig.
Für ein allgemeines Supremum unklassifizierter Größen folgt dies nicht.
Für Log-Volumina, Hüllen, Haar-Normalisierungen und Jacobi-Determinanten
ist hier weder Zugehörigkeit noch Nichtzugehörigkeit bewiesen; ihre
konkrete Darstellung ist die offene Maßpinnungs-Obligation R-E2.
\end{remark}

%=====================================================================
\section{Quellengebundene Anwendung auf die Höhen-Teilkette}\label{sec:sourcebound}

Tabelle~\ref{tab:substitution_audit} trennt das direkte Einsetzergebnis von
Quellenobjekten, die eine zusätzliche Brücke benötigen.
Die Seitenanker beziehen sich auf die zitierten exakten PDFs und behaupten
keine Prüfung ihres gesamten Inhalts.

\begin{table}[htbp]
\centering\small
\caption{Einsetzergebnisse und unbewiesene Bereichs- oder Objektbrücken.}\label{tab:substitution_audit}
\begin{tabularx}{\textwidth}{@{}p{0.26\textwidth}XX@{}}
\toprule
Quellenstelle & Konsequenz im angegebenen Vertrag & Grenze\\
\midrule
II($\tfrac12$) (5.3.3)--(5.3.4), S.~32 & $\beta_y=1$; die gewichtete Produktformel reduziert sich auf die Standardformel & Derselbe Körper $L$ und volles $V_L$\\
\addlinespace
II($\tfrac12$) Def.~8.2.1 / Lem.~8.2.3, S.~51 & $h_y=h_{\mathrm{cl}}$ auf $\mathbb P^n(L)$ & Repräsentantentupel müssen absteigen\\
\addlinespace
II($\tfrac12$) Rem.~8.2.2 / Prop.~8.3.1, S.~51 & Die Quelle erweitert auf $\mathbb P^n(\overline L)$ und behauptet dort nichttriviale Abhängigkeit & Erweiterungs-, Bewertungs- und Basiswechselbrücke hier unbewiesen; durch die $L$-Bereichsidentität allein nicht widerlegt\\
\addlinespace
II($\tfrac12$) \S\S8.6--8.7, S.~54--55 & Auf $L^\times$ ist jeder normalisierte Summand gleich $h_{\mathrm{cl}}$, daher auch das Supremum & Keine Folgerung für andere stabilisierte Objekte\\
\addlinespace
III Rem.~8.9.2(2), S.~90 & Der Höhenvergleich von Punkten in $X(\overline L)$ mittels eines amplen Geradenbündels ist ein gesondertes Quellenobjekt & Brücke zur auditierten projektiven Höhe bleibt offen\\
\addlinespace
IV \S1.5, S.~17 & Die Termvergleichsbeschränkung betrifft untere und obere Log-Volumenvergleiche zwischen $y_0$ und $\varphi(y_0)$ & Kein universelles Verbot aller lokalen Vergleiche; keine $\Geff$-Klassifikation\\
\addlinespace
II($\tfrac12$) Thm.~5.10.1(7) & Koordinatenkerne können bei nichtproportionalen Koeffizienten verschieden sein, während ihre $L^\times$-Rückzüge übereinstimmen & Bedingte Deutung, keine Bestätigung des gesamten Quellensatzes\\
\bottomrule
\end{tabularx}
\end{table}

\emph{Anwendbarkeitsgate H-PLACES.} Der Satz verwendet einen festen
Zahlkörper, alle seine archimedischen und endlichen Stellen und alle
von null verschiedenen Elemente dieses Körpers als Tests.
Eine Quellennotation über nichtarchimedische Stellen
(wie ATS~III \S3.2(5)) darf nicht automatisch mit diesem $V_L$
identifiziert werden. Ebenso verlangen Stellen einer Erweiterung $L'$,
die nur auf $L^\times$ getestet werden, eine explizite Körperabstiegs-
und Stellenbrücke. Keine der beiden Indexidentifikationen folgt aus
einer Koordinatenanalogie.

Die Quelle selbst liefert die punktweise Normalisierung für das Einsetzen
auf dem $L$-Bereich. Damit sind die angegebenen Höhen- und Gradidentitäten
begründet; nicht jedes operative Höhen-, Geradenbündel- oder
Log-Volumenobjekt an anderer Stelle der Serie ist damit identifiziert.

%=====================================================================
\section{Bedingte Konsequenzen (drei typisierte Fälle)}\label{sec:consequences}

Die Fälle betreffen eine geprüfte Größendarstellung und einen ausdrücklich
abgeglichenen Bereich, Körper, eine Stellenfamilie und eine Testmenge.
\begin{description}
\item[Fall A (punktweise Identität (5.3.3)).]\mbox{}\\
Dann ist $\beta_y=1$, und jedes bewiesene Mitglied von $\Geff$ ist gleich
seiner Referenz. Die identifizierte Höhe auf $\mathbb P^n(L)$ und der Grad
auf $L^\times$ sind klassisch und $y$-unabhängig.
Die Anwendung auf eine andere operative Größe verlangt deren
Identifikation und Bereichsbrücke.
\item[Fall B (nur volle gewichtete Produktformel).]\mbox{}\\
Dann ist $\beta_y=c_y\mathbf1$, und für $Q\in\Geff$ gilt
$Q_y=c_yQ_{\mathrm{ref}}$. Der Skalar kann mit $y$ variieren,
die Endgewichte jedoch nicht von Stelle zu Stelle.
Für tatsächliche Standardhöhen und Grade ist dies eine globale
Skalierung auf ihren jeweiligen Bereichen.
\item[Fall C (Zugehörigkeit unklassifiziert).]\mbox{}\\
Für die konkreten Log-Volumina, Hüllen und Maße ist hier weder eine lineare
Darstellung noch ein Hindernis für eine solche bewiesen.
Der Satz beschränkt weiterhin effektive Gewichte, sobald seine
gewichtete Produktformel-Voraussetzung gilt; diese Größenfamilien
klassifiziert er dadurch nicht. R-E2 bleibt offen.
\end{description}

\begin{table}[htbp]
\centering\small
\caption{Typisierte Konsequenzmatrix.}\label{tab:case_verdicts}
\begin{tabularx}{\textwidth}{@{}p{0.16\textwidth}XXX@{}}
\toprule
Fall & Effektive Gewichte & Globale Größen & Herleitungsobligation\\
\midrule
A: punktweise Identität & $\beta_y=1$ & $Q_y=Q_{\mathrm{ref}}$ für bewiesene Mitglieder & Operatives Objekt und Bereichsbrücke identifizieren\\
\addlinespace
B: volle gewichtete Formel & $\beta_y=c_y\mathbf1$ & $Q_y=c_yQ_{\mathrm{ref}}$ für bewiesene Mitglieder & Gewichtete Formel, Indizierung und Faktorisierung prüfen, falls andere Variation gemeint ist\\
\addlinespace
C: unklassifizierte Größe & Bei denselben Satzvoraussetzungen beschränkt & Keine Klassifikation der Größe aus Gewichten allein & Konkrete Darstellung und Maßpinnungsbrücke klären\\
\bottomrule
\end{tabularx}
\end{table}

\emph{Bedingte Konsequenz.} Eine stellenselektive metrische Variation kann
unter der vollen gewichteten Formel nicht von den geprüften linearen
Familien auf dem $L$-Bereich getragen werden; die punktweise Identität
beseitigt sogar ihre globale Skalierung.
Prop.~8.3.1 auf $\mathbb P^n(\overline L)$, der Geradenbündelvergleich
aus ATS~III und der Log-Volumenapparat aus ATS~IV sind gesonderte Ziele.
Ihre operativen Objekte und Transportbrücken müssen geprüft werden,
bevor diese Konsequenz angewandt wird. Die Standard-Produktformel des
festen $b$ ist nichts, dessen Widerlegung diese Fälle von einer Quelle
verlangen. Es folgen weder globaler Kollaps noch Abschluss des
Reparaturprogramms noch Ausschluss anderer Wege.

\emph{Falsifikations- und Revisionskriterien.}
\begin{enumerate}
\item Eine nichtkonstante Effektivgewichtsfamilie, die genau die Produktformel
für denselben Körper und alle Stellen auf jedem $x\in L^\times$ erfüllt,
würde Satz~\ref{thm:rigidity} widersprechen.
Ein Gegenbeispiel zur Faktorisierung muss außerdem den linearen
Darstellungs- und Konvergenzvertrag erfüllen.
\item Der Nachweis eines anderen Bereichs oder einer anderen Definition einer
operativen Quellenhöhe widerlegt einen unbewiesenen Identifikationsschritt,
nicht den Satz. Deshalb bleiben Erweiterungs- und Geradenbündelbrücken
gesondert.
\item Eine künftige Quellenrevision kann neue Objekte oder Brücken liefern.
Sie verlangt eine neue Prüfung dieser Version; sie beweist weder
rückwirkend eine fehlende Brücke noch widerlegt sie eine an die zitierte
Version gebundene Aussage.
\end{enumerate}

%=====================================================================
\section[Externer Programmstatus (nicht Teil der vorliegenden Beweise)]{Externer Programmstatus\texorpdfstring{\\}{ }(nicht Teil der vorliegenden Beweise)}\label{sec:notices}

Teil~5 ist veröffentlicht \cite{part5}; seine Veröffentlichung erfüllt
die verbleibenden mathematischen Brücken nicht.
R6-1 / REP-6-1 ist eine bedingte algebraische
Boxtransport-Rekonstruktion: tatsächliche Haar-Normalisierung,
Tensor-/Indexkompatibilität und der operative Volumenvergleich bleiben
gesonderte Aufgaben. Ein Defekt in einem gewichteten Zwischenmodell
ist für sich kein Gegenbeispiel zu einem Joshi-Theta-Volumensatz.
R6-2 / REP-6-2 setzt V1 und eine explizite $L'\to L$-Stellenbrücke voraus.
R6-3 / REP-6-3 betrifft die Eindeutigkeit normalisierter Daten über
demselben Körper, nicht die Gleichheit aller globalen Koordinatenhyperebenen.
Diese Hinweise beziehen sich auf das Augustprogramm und seine
R18-Korrekturen; keiner wird in den obigen Beweisen verwendet.

Das GNC-Quellenaudit wird in Version 1.2 \cite{gnc}, das
Ledger-Audit aus Teil~4 in Version 0.3 \cite{part4} zitiert.
Das Zertifikatsform-Kriterium ist ein gesondertes Instrument \cite{part2}.
Die Maßpinnungs-Obligation R-E2 aus Teil~6 ist nicht GNC-E2.
Ebenso sind REP-6-1--3 Reparaturpositionen, keine Ledger-F-R-Kennungen;
Publikationsversion 0.3 behauptet keine neue v1-F-Ledger-Abnahme.
Die P5-/Schattentomograph- und Gesamtkettengates bleiben offen.

%=====================================================================
\section{Anspruchsgrenzen}\label{sec:claims}

Satz~\ref{thm:rigidity} und Lemma~\ref{lem:redistribution} betreffen den
expliziten Vertrag für einen festen Körper und alle Stellen.
Lemma~\ref{lem:factorization} verlangt zusätzlich eine geprüfte lineare
Darstellung; Proposition~\ref{prop:exhaustion} betrifft
$\mathbb P^n(L)$ und $L^\times$.
Nur diese geprüften Größen-/Bereichspaare tragen die negative Aussage
über verbleibende stellenselektive Variation.
Die $\overline L$-Erweiterung, Geradenbündelhöhen aus ATS~III und
Log-Volumina aus ATS~IV sind hier nicht mit diesem Vertrag identifiziert.
Behauptet werden weder erschöpfende Abwesenheit eines anderen metrischen
Trägers noch Nichtzugehörigkeit in Fall C, Abschluss des Reparaturprogramms
oder ein Gesamtfehlerverdikt über die Quellen.
Die Koordinatenkerndeutung ist bedingt und bestätigt nicht die vollständige
Nichtkonstanzbehauptung aus Theorem~5.10.1(7) der Quelle.
Es gibt weder Beweis noch Widerlegung von abc oder IUT und keine
mathematische Gesamtabnahme der Serie.

Die R18-Inventur ist die Quelle dieser Korrekturen.
Unabhängige Manuskript- und Quellenprüfungen sind Reviewbelege, kein
maschinengeprüfter Lean-Beweis, kein numerisches Experiment und kein
neues formales Verifikationsergebnis.

\section*{KI-Offenlegung}
KI-Systeme unterstützten Entwurf, Übersetzung und Prüfung dieses
Arbeitspapiers. Die Korrekturversion verwendet die R18-Inventur sowie
unabhängig geprüfte Manuskripte und ausgewählte Quellenstellen.
Der Autor bleibt für den Text und seine Aussagen verantwortlich.
Diese Prüfungen sind keine formale Beweisverifikation und keine
mathematische Gesamtabnahme der ATS-Serie. Es wird weder ein
maschinengeprüfter Beweis noch ein neues numerisches Experiment berichtet.


%=====================================================================
\begin{thebibliography}{99}
\raggedright

\bibitem{artinwhaples}
E.~Artin and G.~Whaples,
\emph{Axiomatic characterization of fields by the product formula for
valuations},
Bull.\ Amer.\ Math.\ Soc.\ \textbf{51} (1945), no.~7, 469--492.
DOI: \href{https://doi.org/10.1090/S0002-9904-1945-08383-9}{10.1090/S0002-9904-1945-08383-9}

\bibitem{neukirch}
J.~Neukirch,
\emph{Algebraic Number Theory},
Grundlehren der mathematischen Wissenschaften \textbf{322},
Springer, Berlin--Heidelberg, 1999.
ISBN 3-540-65399-6. DOI: \href{https://doi.org/10.1007/978-3-662-03983-0}{10.1007/978-3-662-03983-0}

\bibitem{joshiI}
K.~Joshi,
\emph{Construction of Arithmetic Teichmüller Spaces I},
arXiv:2106.11452v4 [math.AG], 2021 (revised 24 February 2025).

\bibitem{joshiII}
K.~Joshi,
\emph{Construction of Arithmetic Teichmüller spaces II: Proof of a local
prototype of Mochizuki's Corollary 3.12},
arXiv:2303.01662v3 [math.AG], 2023 (revised 24 February 2025).

\bibitem{joshiIIhalf}
K.~Joshi,
\emph{Construction of Arithmetic Teichmüller Spaces II$\frac{1}{2}$:
Deformations of Number Fields},
arXiv:2305.10398v12 [math.AG], 2023 (revised 24 February 2025).

\bibitem{joshiIII}
K.~Joshi,
\emph{Construction of Arithmetic Teichmüller Spaces III: A `Rosetta Stone'
and a proof of Mochizuki's Corollary 3.12},
arXiv:2401.13508v4 [math.AG], 2024 (revised 24 February 2025).

\bibitem{joshiIV}
K.~Joshi,
\emph{Construction of Arithmetic Teichmüller Spaces IV: Proof of the
abc-conjecture},
arXiv:2403.10430v2 [math.AG], 2024 (revised 24 February 2025).

\bibitem{mochizukiIUT}
S.~Mochizuki,
\emph{Inter-universal Teichmüller Theory I--IV},
Publ.\ Res.\ Inst.\ Math.\ Sci.\ \textbf{57} (2021), no.~1/2, 3--723.

\bibitem{part1}
L.~Geiger,
\emph{The Comparison in IUTchIII, Corollary 3.12: From Attempted Repair to Structural Diagnosis},
Zenodo, 2026.
Concept DOI: \href{https://doi.org/10.5281/zenodo.19960781}{10.5281/zenodo.19960781}.
Zitierte Fassung v6.4, 2026-10-01.
DOI: \href{https://doi.org/10.5281/zenodo.23081604}{10.5281/zenodo.23081604}.

\bibitem{part3}
L.~Geiger,
\emph{The Bridge Certificate: A Line-Level Audit Instrument for Inter-Theoretic Transfer Arguments, with Four Applications to the abc Literature},
Zenodo, 2026.
Concept DOI: \href{https://doi.org/10.5281/zenodo.21925803}{10.5281/zenodo.21925803}.
Zitierte Fassung 0.3, 2026-10-01.
DOI: \href{https://doi.org/10.5281/zenodo.23090371}{10.5281/zenodo.23090371}.

\bibitem{part4}
L.~Geiger,
\emph{The Bridge Certificate Applied to the Joshi ATS Series: A Ledger Audit (DRAFT)},
Zenodo, 2026.
Concept DOI: \href{https://doi.org/10.5281/zenodo.21951155}{10.5281/zenodo.21951155}.
Zitierte Fassung 0.3, 2026-10-02.
DOI: \href{https://doi.org/10.5281/zenodo.23092802}{10.5281/zenodo.23092802}.

\bibitem{gnc}
L.~Geiger,
\emph{Global Normalization in Joshi's Arithmetic Teichmüller Theory: What the Identification Carries --- and What Carries the Theorem (DRAFT)},
Zenodo, 2026.
Concept DOI: \href{https://doi.org/10.5281/zenodo.21924253}{10.5281/zenodo.21924253}.
Zitierte Fassung 1.2, 2026-10-02.
DOI: \href{https://doi.org/10.5281/zenodo.23092166}{10.5281/zenodo.23092166}.

\bibitem{landscapeA}
L.~Geiger,
\emph{The abc Landscape: Obstructions, Failed Routes, and HCT Diagnostics (DRAFT)},
Zenodo, 2026.
Concept DOI: \href{https://doi.org/10.5281/zenodo.21916900}{10.5281/zenodo.21916900}.
Zitierte Fassung 1.1, 2026-08-13.
DOI: \href{https://doi.org/10.5281/zenodo.21924338}{10.5281/zenodo.21924338}.

\bibitem{part2}
L.~Geiger,
\emph{The Bridge-Contract Criterion: A Form Audit for Transport Arguments, with an Instantiation in the abc Literature},
Zenodo, 2026.
Concept DOI: \href{https://doi.org/10.5281/zenodo.21960476}{10.5281/zenodo.21960476}.
Zitierte Fassung 0.2, 2026-08-16.
DOI: \href{https://doi.org/10.5281/zenodo.21960477}{10.5281/zenodo.21960477}.

\bibitem{part5}
L.~Geiger,
\emph{The Joshi ATS Series under the Bridge-Certificate Standard: A Targeted Substance Audit and an All-Place Rigidity Classification of Effective Normalization Weights},
Zenodo, 2026.
Concept DOI: \href{https://doi.org/10.5281/zenodo.21956398}{10.5281/zenodo.21956398}.
Zitierte Fassung 0.2, 2026-08-15.
DOI: \href{https://doi.org/10.5281/zenodo.21956399}{10.5281/zenodo.21956399}.

\end{thebibliography}

\end{document}
